% Vietnamese translation for OI Public Library
% Source-PDF: 比赛 CONTEST/2021“MINIEYE杯”中国大学生算法设计超级联赛/2021“MINIEYE杯”中国大学生算法设计超级联赛（7）/2021“MINIEYE杯”中国大学生算法设计超级联赛（7）-题集.pdf
\documentclass[11pt]{article}
\usepackage{oipl-vn}

\title{Đề bài trận thứ bảy HDU Multi-University 2021}
\author{}
\date{}

\begin{document}
\maketitle

\origtitle{Super League of Chinese College Students Algorithm Design 2021, Contest 7}
\sourcepath{contest/hdu-2021-minieye-07-problems.pdf}

\section{Problem 1001. Fall with Fake Problem}

\paragraph{Tệp vào:} standard input
\paragraph{Tệp ra:} standard output

Một cậu bé có ID là smzzl rất thích giải bài trên Codeforces. Một ngày nọ,
cậu đang giải bài toán sau:

Cho một chuỗi \(S\), hãy tìm chuỗi \(T\) nhỏ nhất theo thứ tự từ điển, sao cho
thứ tự từ điển của \(T\) không nhỏ hơn \(S\), số lần xuất hiện của mỗi chữ cái
trong \(T\) đều là bội của \(k\), và \(|S|=|T|\).

``Bài này dễ quá! Mình muốn một bài khó hơn!''

Vì vậy, cậu sửa bài toán. Cậu không còn yêu cầu số lần xuất hiện của mỗi chữ
cái trong \(T\) là bội của \(k\). Thay vào đó, cậu muốn chúng là một ước của
\(k\), hoặc bằng \(0\).

Bài toán thật sự trở nên khó hơn, nhưng cũng không khó hơn quá nhiều. Fall,
đồng đội của smzzl, đã giải được ngay lập tức. Bây giờ Fall muốn bạn giải bài
này để cậu có thể kiểm tra đáp án của mình.

\subsection*{Dữ liệu vào}

Dữ liệu vào gồm nhiều bộ dữ liệu.

Dòng đầu tiên chứa một số nguyên \(T\) (\(1\le T\le 32\)) là số bộ dữ liệu.

Với mỗi bộ dữ liệu có hai dòng.

Dòng đầu tiên chứa hai số nguyên \(n,k\) (\(1\le n,k\le 10^5\)), lần lượt là
độ dài của \(S\) và tham số đã nêu ở trên.

Dòng thứ hai chứa chuỗi \(S\), chỉ gồm các chữ cái tiếng Anh thường.

Bảo đảm rằng tổng \(|S|\) trên tất cả các bộ dữ liệu không vượt quá \(10^6\).

\subsection*{Dữ liệu ra}

Với mỗi bộ dữ liệu, in một chuỗi \(T\) trên một dòng, là đáp án của bạn. Nếu
không tồn tại chuỗi như vậy, in \(-1\) thay vào đó.

\subsection*{Ví dụ}

\paragraph{standard input}
\begin{verbatim}
2
5 6
aaaaa
10 18
abcabcabcd
\end{verbatim}

\paragraph{standard output}
\begin{verbatim}
aaabb
abcabcabcd
\end{verbatim}

\section{Problem 1002. Fall with Soldiers}

\paragraph{Tệp vào:} standard input
\paragraph{Tệp ra:} standard output

Fall đang chơi một trò chơi về chiến tranh.

Có \(n\) người lính đứng thành một hàng. Một số thuộc về người chơi, một số
thuộc về kẻ địch. Do có gián điệp, phe của một vài người lính chưa được biết.
Bây giờ, với vai trò pháo binh, người chơi có thể thắng trò chơi theo các
luật sau sau khi xác định phe của từng người lính:

Bước 1: Chọn một người lính thuộc về người chơi. Người lính được chọn phải có
đúng hai người lính đứng cạnh, nghĩa là người đó không được ở đầu hoặc cuối
hàng.

Bước 2: Giết hai người lính đứng cạnh người được chọn. Người được chọn vẫn còn
sống.

Bước 3: Lặp lại bước 1 và bước 2 cho đến khi tất cả những người lính còn lại
đều thuộc cùng một phe, tức là phe người chơi hoặc phe kẻ địch.

Người chơi thắng nếu tất cả những người lính còn lại đều thuộc về người chơi.
Fall, trong vai người chơi, muốn biết có bao nhiêu cách xác định phe của từng
người lính để cậu có thể thắng.

Tuy nhiên, trạng thái của lính có thể thay đổi, thành quân của người chơi,
kẻ địch, hoặc chưa biết. Bạn cần tính đáp án sau mỗi lần thay đổi.

\subsection*{Dữ liệu vào}

Dữ liệu vào gồm nhiều bộ dữ liệu.

Dòng đầu tiên chứa một số nguyên \(T\) (\(1\le T\le 11\)) là số bộ dữ liệu.

Với mỗi bộ dữ liệu:

Dòng đầu tiên chứa hai số nguyên \(n,q\)
(\(1\le n,q\le 2\times 10^5\), \(n\) là số lẻ), lần lượt là số người lính và
số thao tác.

Dòng thứ hai chứa một chuỗi \(s\) độ dài \(n\), biểu thị trạng thái ban đầu
của các người lính. Nếu \(s_i\) là `1', người lính đó thuộc về bạn. Nếu
\(s_i\) là `0', người lính đó thuộc về kẻ địch. Nếu \(s_i\) là `?', phe của
người lính đó chưa biết.

Trong \(q\) dòng tiếp theo, mỗi dòng chứa một số nguyên \(p\) (\(1\le p\le |s|\))
và một ký tự \(c\), nghĩa là đổi \(s_p\) thành \(c\).

Bảo đảm rằng tổng \(q\) trên tất cả các bộ dữ liệu không vượt quá \(10^6\), và
\(s_i,c\in\{\text{`0'},\text{`1'},\text{`?'}\}\).

\subsection*{Dữ liệu ra}

Với mỗi bộ dữ liệu, in đáp án ban đầu trên một dòng. Sau đó in \(q\) đáp án
trên \(q\) dòng, là đáp án sau từng lần thay đổi. Tất cả đáp án cần được in
theo modulo \(10^9+7\).

\subsection*{Ví dụ}

\paragraph{standard input}
\begin{verbatim}
6
5 1
01000
1 1
5 1
00000
3 1
15 4
1100????????111
8 ?
1 0
4 1
11 ?
15 4
0????000?0???0?
6 0
10 0
1 ?
10 ?
15 2
???????0?0???0?
4 0
1 0
15 2
00000000?0???0?
4 ?
4 ?
\end{verbatim}

\paragraph{standard output}
\begin{verbatim}
0
0
0
1
247
247
247
253
253
368
368
368
736
1664
3616
1760
880
0
24
24
\end{verbatim}

\section{Problem 1003. Fall with Trees}

\paragraph{Tệp vào:} standard input
\paragraph{Tệp ra:} standard output

Fall muốn vẽ một cây nhị phân hoàn hảo.

Trước hết, ta quy ước rằng tất cả các nút trong cây có cùng độ sâu cũng có
cùng tọa độ \(y\) trên mặt phẳng. Gọi các nút có cùng độ sâu là các nút cùng
tầng. Khi đó cây nhị phân hoàn hảo có bốn tính chất:

\begin{itemize}
  \item Nó là một cây nhị phân đầy đủ.
  \item Hiệu tọa độ \(y\) của hai nút ở hai tầng kề nhau là một hằng số.
  \item Hiệu tọa độ \(x\) của hai nút kề nhau trên cùng một tầng là một hằng
  số.
  \item Tọa độ \(x\) của mỗi nút bằng trung bình cộng tọa độ \(x\) của hai con
  của nó.
\end{itemize}

Fall đã vẽ nút gốc và hai con trái, phải của nút gốc trong cây nhị phân này.
Bây giờ Fall định vẽ tổng cộng \(k\) tầng, sau đó cắt cây nhị phân ra và dán
lên tường. Vì vậy cậu muốn biết diện tích bao lồi của tất cả các nút trong
cây nhị phân hoàn hảo này.

Trong hình minh họa của ví dụ đầu tiên, diện tích là \(S_{ABC}+S_{BCGD}=14\).

\subsection*{Dữ liệu vào}

Dữ liệu vào gồm nhiều bộ dữ liệu.

Dòng đầu tiên chứa một số nguyên \(T\) (\(T\le 2\times 10^5\)) là số bộ dữ
liệu.

Với mỗi bộ dữ liệu:

Dòng đầu tiên chứa một số nguyên \(k\) (\(2\le k\le 10^4\)).

Dòng thứ hai chứa sáu số nguyên
\[
  x_{\text{root}},y_{\text{root}},x_{\text{lson}},y_{\text{lson}},
  x_{\text{rson}},y_{\text{rson}}\in[-10^4,10^4],
\]
biểu thị tọa độ của nút gốc và hai con của nó.

Bảo đảm rằng tất cả tọa độ thỏa mãn các điều kiện của đề bài, nghĩa là:

\begin{itemize}
  \item \(x_{\text{lson}}+x_{\text{rson}}=2\times x_{\text{root}}\).
  \item \(y_{\text{lson}}=y_{\text{rson}}\).
  \item \(y_{\text{root}}>y_{\text{lson}}\), \(x_{\text{lson}}<x_{\text{rson}}\).
\end{itemize}

\subsection*{Dữ liệu ra}

Với mỗi bộ dữ liệu, in một số thực biểu thị đáp án, với ba chữ số sau dấu
thập phân.

\subsection*{Ví dụ}

\paragraph{standard input}
\begin{verbatim}
3
3
0 0 -2 -2 2 -2
4
0 0 -4 -2 4 -2
10000
0 0 -10000 -10000 10000 -10000
\end{verbatim}

\paragraph{standard output}
\begin{verbatim}
14.000
54.000
3999000000000.000
\end{verbatim}

\section{Problem 1004. Link with Balls}

\paragraph{Tệp vào:} standard input
\paragraph{Tệp ra:} standard output

Trong nhà máy của Ball Illusion Technology (BIT) có rất nhiều quả bóng. Link,
một cậu bé, đến đó để lấy vài quả bóng, nhưng đột nhiên cậu nhận ra rằng có
quá nhiều cách lấy bóng.

Có \(2n\) chiếc xô trong nhà máy. Link có thể lấy \(kx\) quả bóng từ chiếc xô
thứ \(2x-1\), trong đó \(k\) là một số nguyên không âm. Cậu cũng có thể lấy
nhiều nhất \(x\) quả bóng từ chiếc xô thứ \(2x\).

Link muốn lấy \(m\) quả bóng, và cậu tự hỏi có bao nhiêu cách lấy ra đúng
\(m\) quả bóng. Trong khi Link đang tính đáp án, cậu cũng muốn bạn tính nó, và
bạn cần in đáp án theo modulo \(10^9+7\).

\subsection*{Dữ liệu vào}

Dữ liệu vào gồm nhiều bộ dữ liệu.

Dòng đầu tiên chứa một số nguyên \(T\) (\(1\le T\le 10^5\)) là số bộ dữ liệu.

Mỗi bộ dữ liệu chứa hai số nguyên \(n\) và \(m\)
(\(1\le n,m\le 10^6\)).

\subsection*{Dữ liệu ra}

Với mỗi bộ dữ liệu, in đáp án theo modulo \(10^9+7\) trên một dòng.

\subsection*{Ví dụ}

\paragraph{standard input}
\begin{verbatim}
4
1 1
2 2
3 3
1000000 1000000
\end{verbatim}

\paragraph{standard output}
\begin{verbatim}
2
6
20
192151600
\end{verbatim}

\section{Problem 1005. Link with EQ}

\paragraph{Tệp vào:} standard input
\paragraph{Tệp ra:} standard output

Ta thường nói rằng các chi tiết có thể phản ánh mức EQ của một người. Ví dụ,
ở BIT, khi sinh viên học trong phòng học, chẳng hạn A105 trong tòa nhà giảng
dạy chung, hoặc học ở các bàn dài trong thư viện, sẽ thường rất ngượng ngùng
và khó chịu, đôi khi còn có mùi khó chịu, nếu một sinh viên khác đột nhiên
ngồi ngay cạnh họ. Ta gọi đó là hành vi EQ thấp.

Bây giờ, có một chiếc bàn dài gồm \(n\) chỗ ngồi, và sinh viên BIT sẽ chọn chỗ
ngồi theo cách sau để tránh hành vi EQ thấp nhiều nhất có thể.

\begin{itemize}
  \item Nếu cả bàn đang trống, chọn ngẫu nhiên một vị trí để ngồi.
  \item Ngược lại, chọn một vị trí sao cho khoảng cách ngắn nhất đến các sinh
  viên khác là lớn nhất. Nếu có nhiều vị trí như vậy, chọn ngẫu nhiên một
  trong số đó với xác suất bằng nhau.
\end{itemize}

Khi khoảng cách xa nhất bằng \(1\), tức là dù sao sinh viên cũng sẽ ngồi cạnh
một bạn học, sinh viên sẽ ngồi ở một bàn khác để tránh hành vi EQ thấp, và
khi đó ta gọi chiếc bàn này là đã đầy.

Bây giờ Link muốn biết kỳ vọng số người đã ngồi ở bàn khi bàn trở thành đầy.

\subsection*{Dữ liệu vào}

Dữ liệu vào gồm nhiều bộ dữ liệu.

Dòng đầu tiên chứa một số nguyên \(T\) (\(1\le T\le 10^5\)) là số bộ dữ liệu.

Với mỗi bộ dữ liệu, chỉ có một số nguyên \(n\) (\(1\le n\le 10^6\)) trên một
dòng, là độ dài của bàn.

\subsection*{Dữ liệu ra}

Với mỗi bộ dữ liệu, in đáp án theo modulo \(10^9+7\).

Bạn có thể tìm cách lấy modulo của một số hữu tỉ trong Problem 1006.

\subsection*{Ví dụ}

\paragraph{standard input}
\begin{verbatim}
5
1
2
3
4
2021
\end{verbatim}

\paragraph{standard output}
\begin{verbatim}
1
1
666666673
2
420089906
\end{verbatim}

\section{Problem 1006. Link with Grenade}

\paragraph{Tệp vào:} standard input
\paragraph{Tệp ra:} standard output

``Chúng ta đang lãng phí thời gian, đi thôi!''

``Khủng bố thắng.''

``Tập hợp đội hình, đi thôi!''

``Khủng bố thắng.''

\(\ldots\)

Sau khi bị đánh bại nhiều lần, cậu tức giận đến mức quyết định gian lận. Cậu
dùng vài thủ thuật để khiến mình bay thật cao, để không ai có thể làm hại cậu.
Cậu bắt đầu ném lựu đạn, nhưng đột nhiên nhận ra mình chỉ có \(1\) HP, nghĩa
là cậu sẽ chết nếu lựu đạn làm cậu bị thương. Bây giờ Link muốn biết xác suất
cậu sẽ chết.

Một cách hình thức, có một quả lựu đạn với vận tốc ban đầu \(v_0\) m/s và
thời gian nổ \(t\) giây, được ném theo một hướng ngẫu nhiên. Khi lựu đạn nổ,
nó sẽ làm bị thương bất kỳ ai có khoảng cách đến nó không quá \(r\). Giả sử
lựu đạn không chạm vào bất cứ thứ gì trước khi nổ, và người ném không di
chuyển sau khi ném lựu đạn, bạn cần in xác suất người đó sống sót.

Để kiểm tra đáp án chính xác, bạn cần in đáp án theo modulo \(10^9+7\).

Ghi chú 1: Đây là một trò chơi 3D. Trong trò chơi này, gia tốc trọng trường là
\(10\ \text{m/s}^2\).

Ghi chú 2: Có thể chứng minh rằng đáp án luôn biểu diễn được dưới dạng
\(\frac{p}{q}\), trong đó ước chung lớn nhất của \(p\) và \(q\) là \(1\). Khi
in đáp án theo modulo \(M\), bạn cần in một số nguyên \(x\) sao cho
\(0\le x<M\) và \(x\cdot q\equiv p\pmod M\).

\subsection*{Dữ liệu vào}

Dữ liệu vào gồm nhiều bộ dữ liệu.

Dòng đầu tiên chứa một số nguyên \(T\) (\(1\le T\le 10^5\)) là số bộ dữ liệu.

Mỗi bộ dữ liệu chứa ba số nguyên \(t_0\), \(v_0\) và \(R\)
(\(1\le t_0,v_0,R\le 100\)).

\subsection*{Dữ liệu ra}

Với mỗi bộ dữ liệu, in đáp án theo modulo \(10^9+7\) trên một dòng.

\subsection*{Ví dụ}

\paragraph{standard input}
\begin{verbatim}
5
1 10 10
1 16 10
1 15 10
1 4 10
1 5 10
\end{verbatim}

\paragraph{standard output}
\begin{verbatim}
625000005
1
1
0
0
\end{verbatim}

\section{Problem 1007. Link with Limit}

\paragraph{Tệp vào:} standard input
\paragraph{Tệp ra:} standard output

Link có một hàm \(f(x)\), trong đó \(x\) và \(f(x)\) đều là các số nguyên
trong \([1,n]\).

Gọi \(f_n(x)=f(f_{n-1}(x))\) và \(f_1(x)=f(x)\). Cậu định nghĩa sức mạnh của
một số \(x\) là:
\[
  g(x)=\lim_{n\to +\infty}\frac{1}{n}\sum_{i=1}^{n} f_i(x).
\]

Cậu muốn biết liệu mọi \(x\in[1,n]\) có cùng sức mạnh hay không.

\subsection*{Dữ liệu vào}

Dữ liệu vào gồm nhiều bộ dữ liệu.

Dòng đầu tiên chứa một số nguyên \(T\) (\(1\le T\le 100\)) là số bộ dữ liệu.

Với mỗi bộ dữ liệu:

Dòng đầu tiên chứa một số nguyên \(n\) (\(1\le n\le 10^5\)).

Dòng thứ hai chứa \(n\) số nguyên, số nguyên thứ \(i\) cho biết giá trị của
\(f(i)\) (\(1\le f(i)\le n\)).

Bảo đảm rằng tổng \(n\) trên tất cả các bộ dữ liệu không vượt quá \(10^6\).

\subsection*{Dữ liệu ra}

Với mỗi bộ dữ liệu, in `YES' nếu mọi \(x\) có cùng sức mạnh. Ngược lại, in
`NO'.

\subsection*{Ví dụ}

\paragraph{standard input}
\begin{verbatim}
2
2
1 2
2
1 1
\end{verbatim}

\paragraph{standard output}
\begin{verbatim}
NO
YES
\end{verbatim}

\section{Problem 1008. Smzzl with Greedy Snake}

\paragraph{Tệp vào:} standard input
\paragraph{Tệp ra:} standard output

Smzzl sắp tạo một AI cho trò Rắn Săn Mồi. Trò chơi diễn ra trên mặt phẳng
\(xOy\) và không có vật cản nào trên mặt phẳng.

Trong trò chơi này, con rắn cần \(1\) đơn vị thời gian để đi thẳng thêm một
đơn vị độ dài. Nó cũng cần \(1\) đơn vị thời gian để xoay \(90\) độ. Con rắn
phải xoay trong trọn một đơn vị thời gian. Ban đầu có một thức ăn trên bản đồ.
Sau khi rắn ăn mỗi thức ăn, thức ăn tiếp theo xuất hiện.

Smzzl tất nhiên muốn con rắn ăn thức ăn nhanh nhất có thể, vì vậy cậu cần
tối thiểu hóa thời điểm con rắn ăn từng thức ăn. Hãy in một dãy thao tác hợp
lệ.

\subsection*{Dữ liệu vào}

Dữ liệu vào gồm nhiều bộ dữ liệu.

Dòng đầu tiên chứa một số nguyên \(T\) (\(1\le T\le 200\)) là số bộ dữ liệu.

Với mỗi bộ dữ liệu:

Dòng đầu tiên chứa ba số nguyên \(x,y,d\)
(\(|x|,|y|\le 10^4,\ 0\le d\le 3\)). Con rắn bắt đầu tại \((x,y)\). \(d\) cho
biết hướng của đầu rắn: \(0\) là hướng \(y+\), \(1\) là hướng \(x+\), \(2\) là
hướng \(y-\), và \(3\) là hướng \(x-\).

Dòng thứ hai chứa một số nguyên \(n\) (\(1\le n\le 10^5\)), là số thức ăn.

Trong \(n\) dòng tiếp theo, mỗi dòng chứa hai số nguyên \(x,y\)
(\(|x|,|y|\le 10^4\)), nghĩa là thức ăn tiếp theo xuất hiện tại \((x,y)\).

Bảo đảm rằng mọi đường thẳng nối hai thức ăn xuất hiện liên tiếp không song
song với trục \(x\) hoặc trục \(y\).

\subsection*{Dữ liệu ra}

Với mỗi bộ dữ liệu, in dãy thao tác ngắn nhất. In `f' cho đi thẳng, `c' cho
xoay theo chiều kim đồng hồ, và `u' cho xoay ngược chiều kim đồng hồ. Mỗi thao
tác kéo dài một đơn vị thời gian.

Có thể chứng minh rằng chỉ có một dãy thao tác thỏa mãn yêu cầu.

Bảo đảm rằng tổng độ dài dữ liệu ra không vượt quá \(2\times 10^6\).

\subsection*{Ví dụ}

\paragraph{standard input}
\begin{verbatim}
2
0 0 0
2
-1 -1
1 1
0 0 2
5
-1 2
2 4
3 -5
4 -2
5 0
\end{verbatim}

\paragraph{standard output}
\begin{verbatim}
ufufuffuff
cfcffffcffffcfffffffffufufffffcf
\end{verbatim}

\section{Problem 1009. Smzzl with Safe Zone}

\paragraph{Tệp vào:} standard input
\paragraph{Tệp ra:} standard output

Một cậu bé có ID là smzzl đang chơi một trò chơi tên là Zig-Zag-Land (ZZL).
Đây là một trò chơi battle royale 2D. Khi trò chơi bắt đầu, bạn sẽ xuất hiện
tại một vị trí ngẫu nhiên trên bản đồ, và mục tiêu của bạn là di chuyển vào
vùng an toàn.

Cho biên của bản đồ và biên của vùng an toàn, cậu muốn biết khoảng cách ngắn
nhất giữa điểm xuất hiện và vùng an toàn trong trường hợp xấu nhất.

Có thể chứng minh rằng bình phương của đáp án luôn là số hữu tỉ. Để kiểm tra
đáp án chính xác, bạn cần in bình phương đáp án theo modulo \(10^9+7\). Bạn
có thể tìm cách lấy modulo của một số hữu tỉ trong Problem 1006.

\subsection*{Dữ liệu vào}

Dữ liệu vào gồm nhiều bộ dữ liệu.

Dòng đầu tiên chứa một số nguyên \(T\) (\(1\le T\le 10^4\)) là số bộ dữ liệu.

Với mỗi bộ dữ liệu:

Dòng đầu tiên chứa một số nguyên \(m\) (\(3\le m\le 10^5\)), là số đỉnh của
vùng an toàn.

Trong \(m\) dòng tiếp theo, mỗi dòng chứa hai số nguyên \(x_i,y_i\)
(\(|x_i|,|y_i|\le 10^7\)), là một đỉnh của vùng an toàn.

Dòng tiếp theo chứa một số nguyên \(n\) (\(3\le n\le 10^5\)), là số đỉnh của
bản đồ trò chơi.

Trong \(n\) dòng tiếp theo, mỗi dòng chứa hai số nguyên \(x_i,y_i\)
(\(|x_i|,|y_i|\le 10^7\)), là một đỉnh của bản đồ trò chơi.

Bảo đảm rằng:

\begin{itemize}
  \item Vùng an toàn là một tập con của bản đồ trò chơi.
  \item Vùng an toàn và bản đồ trò chơi đều là các đa giác lồi.
  \item Tất cả các đỉnh được cho theo thứ tự ngược chiều kim đồng hồ.
  \item Tổng \(n\) trên tất cả các bộ dữ liệu không vượt quá \(10^6\).
  \item Tổng \(m\) trên tất cả các bộ dữ liệu không vượt quá \(10^6\).
\end{itemize}

Bạn có thể cần dùng phương thức nhập nhanh hơn.

\subsection*{Dữ liệu ra}

Với mỗi bộ dữ liệu, in bình phương đáp án theo modulo \(10^9+7\).

\subsection*{Ví dụ}

\paragraph{standard input}
\begin{verbatim}
1
4
0 0
1 0
1 1
0 1
4
-1 -1
2 -1
2 2
-1 2
\end{verbatim}

\paragraph{standard output}
\begin{verbatim}
2
\end{verbatim}

\section{Problem 1010. Smzzl with Tropical Taste}

\paragraph{Tệp vào:} standard input
\paragraph{Tệp ra:} standard output

Một cậu bé có ID là smzzl rất thích uống Black Ice Tea (BIT), đặc biệt là vị
nhiệt đới.

Chủ cửa hàng biết điều đó, và cô chuẩn bị một bể bơi để có đủ không gian chứa
BIT. Khi cậu bé biết rằng chủ cửa hàng đã chuẩn bị một bể bơi, cậu bắt đầu
uống trong bể. Trong khi chủ cửa hàng đang đổ BIT vào bể bơi, cậu bé cũng
đang uống nó.

Giả sử trong bể bơi có \(V\) lít BIT, tốc độ chủ cửa hàng đổ BIT là \(qV\)
lít mỗi giây, và tốc độ cậu bé uống BIT là \(pV\) lít mỗi giây. Lưu ý rằng
\(V\) thay đổi theo thời gian.

Bây giờ, chủ cửa hàng muốn biết mệnh đề sau có đúng hay không:

Với mọi \(G>0\), tồn tại \(T>0\) sao cho với mọi \(t>T\), cậu bé uống nhiều
hơn \(G\) lít BIT sau \(t\) giây.

\subsection*{Dữ liệu vào}

Dữ liệu vào gồm nhiều bộ dữ liệu.

Dòng đầu tiên chứa một số nguyên \(T\) (\(1\le T\le 100\)) là số bộ dữ liệu.

Với mỗi bộ dữ liệu, có hai số thập phân \(p,q\)
(\(0<p,q<10^4\), chúng có nhiều nhất \(4\) chữ số sau dấu thập phân) trên một
dòng. \(p,q\) là các giá trị đã nêu trong đề bài. Bạn có thể xem thể tích ban
đầu của BIT là \(1\) lít, và bể bơi có thể chứa vô hạn lít BIT.

\subsection*{Dữ liệu ra}

Với mỗi bộ dữ liệu, in một câu trên một dòng.

Nếu mệnh đề là đúng, in `N0 M0R3 BL4CK 1CE TEA!'.

Nếu mệnh đề là sai, in `ENJ0Y YOURS3LF!'.

\subsection*{Ví dụ}

\paragraph{standard input}
\begin{verbatim}
2
1.1 2.2
2.05 1.4
\end{verbatim}

\paragraph{standard output}
\begin{verbatim}
N0 M0R3 BL4CK 1CE TEA!
ENJ0Y YOURS3LF!
\end{verbatim}

\section{Problem 1011. Yiwen with Formula}

\paragraph{Tệp vào:} standard input
\paragraph{Tệp ra:} standard output

Cho một mảng \(a\) độ dài \(n\). Với mọi mảng \(b\) thỏa mãn
\(1\le b_i\le n\) và \(b_1<b_2<\cdots<b_k\), trong đó \(k\) là độ dài của
\(b\) và \(k\ge 1\), hãy tính:
\[
  \prod_{b_1<b_2<\cdots<b_k}\left(a_{b_1}+a_{b_2}+\cdots+a_{b_k}\right).
\]

\subsection*{Dữ liệu vào}

Dữ liệu vào gồm nhiều bộ dữ liệu.

Dòng đầu tiên chứa một số nguyên \(T\) (\(1\le T\le 10\)) là số bộ dữ liệu.

Với mỗi bộ dữ liệu:

Dòng đầu tiên chứa một số nguyên \(n\) (\(1\le n\le 10^5\)), là độ dài của
\(a\).

Dòng thứ hai chứa \(n\) số nguyên \(a_i\) (\(0\le a_i\le 10^5\)), là mảng
\(a\).

Bảo đảm rằng:

\begin{itemize}
  \item Tổng \(n\) trên tất cả các bộ dữ liệu không vượt quá
  \(2.5\times 10^5\).
  \item Tổng \(a_i\) trong một bộ dữ liệu không vượt quá \(10^5\).
  \item Tổng \(a_i\) trên tất cả các bộ dữ liệu không vượt quá \(4\times 10^5\).
\end{itemize}

\subsection*{Dữ liệu ra}

Với mỗi bộ dữ liệu, in kết quả theo modulo \(998244353\) trên một dòng.

\subsection*{Ví dụ}

\paragraph{standard input}
\begin{verbatim}
3
2
1 1
3
1 1 2
5
4 6 9 1 5
\end{verbatim}

\paragraph{standard output}
\begin{verbatim}
2
144
417630946
\end{verbatim}

\section{Problem 1012. Yiwen with Sqc}

\paragraph{Tệp vào:} standard input
\paragraph{Tệp ra:} standard output

Gần đây, Yiwen đang học hợp ngữ.

Vì Yiwen là một thiên tài khoa học máy tính, cậu đề xuất một kỹ thuật cực kỳ
có giá trị gọi là Basic Instruction Technology. Dùng kỹ thuật này, Yiwen có
thể tạo ra rất nhiều lệnh khác nhau. Một vài lệnh trong số đó có thể dùng để
tính một trị riêng đặc biệt của chuỗi, chẳng hạn như sqc.

Với một chuỗi \(s\), gọi \(s_i\) là ký tự thứ \(i\) của nó.

Lệnh sqc có \(4\) toán hạng \(s,l,r,c\), trong đó \(s\) là một chuỗi,
\(1\le l\le r\le n\), và \(c\) là một ký tự được biểu diễn bằng giá trị ASCII.
\(\operatorname{sqc}(s,l,r,c)\) trả về số ký tự \(c\) trong chuỗi con
\(s_l\cdots s_r\).

Yiwen muốn bạn tính biểu thức sau với chuỗi \(s\) đã cho:
\[
  \sum_{c=97}^{122}\sum_{i=1}^{n}\sum_{j=i}^{n}
  \operatorname{sqc}(s,i,j,c)^2.
\]

Lưu ý rằng giá trị ASCII của `a' là \(97\) và giá trị ASCII của `z' là \(122\).

Đáp án có thể rất lớn, nên chỉ cần tính đáp án theo modulo \(998244353\).

\subsection*{Dữ liệu vào}

Dữ liệu vào gồm nhiều bộ dữ liệu.

Dòng đầu tiên chứa một số nguyên \(T\) (\(1\le T\le 100\)) là số bộ dữ liệu.

Với mỗi bộ dữ liệu, dòng duy nhất chứa một chuỗi \(s\) (\(|s|\le 10^5\)) gồm
các chữ cái tiếng Anh thường.

Bảo đảm rằng tổng \(|s|\) trên tất cả các bộ dữ liệu không vượt quá
\(5\times 10^6\).

\subsection*{Dữ liệu ra}

Với mỗi bộ dữ liệu, in đáp án theo modulo \(998244353\) trên một dòng.

\subsection*{Ví dụ}

\paragraph{standard input}
\begin{verbatim}
5
ababa
a
abcdefga
afidhehfhd
ggggggggggg
\end{verbatim}

\paragraph{standard output}
\begin{verbatim}
57
1
122
328
1716
\end{verbatim}

\end{document}
